\begin{abstract}
Machine translation (MT) has become an essential tool for software localization, yet the comparative performance of different MT systems on technical documentation and user interface (UI) texts remains understudied, particularly for the English-Turkish language pair. This paper presents a comprehensive evaluation of four major MT tools: Google Translate, DeepL, Microsoft Translator, and Amazon Translate. We compiled a dataset of 50,000+ parallel segments from open-source software localization files, technical documentation, and UI strings, categorized into technical documentation (25,000 segments), UI strings (15,000 segments), and error messages (10,000 segments). Our experiments on 10,000 test segments using automatic metrics (BLEU, METEOR, TER, chrF++) reveal significant performance differences across tools and text categories. Statistical analysis using paired t-tests and ANOVA confirms these differences are highly significant (p < 0.001). We found that [Tool X] achieved the highest average BLEU score of [X.XX] on technical documentation, while [Tool Y] performed best on UI strings with [Y.YY]. Text length significantly impacts quality, with medium-length texts (50-200 characters) yielding optimal results. Placeholder preservation rates varied from [X]\% to [Y]\%, critical for software functionality. We provide practical recommendations for developers and localization teams, considering both quality and cost trade-offs. Our open-source dataset, evaluation framework, and interactive web application are publicly available to facilitate future research and practical applications.
\end{abstract}
